%% kapa: exam=1 round=23 year=2012 period=1 form=A subject=economics subject_ko=경제원론 seq=17 number=57
%% source: Q-Net 감정평가사 기출문제 (cst00309) artlSeq=5052578 file=raw/1008/5052578/2012년 제23회 감정평가사 제1차시험 기출문제(1교시 A형).pdf
%% conversion: kordoc 4.24.51 -> tools/gen1.py (auto); review=proofread-sonnet; flags=none
\begin{question}{57}
한 시장에 두 기업 A, B가 존재한다. 각 기업은 두 가지 생산 전략 L, H 중 하나를 선택할 수 있다. 두 기업의 생산 전략 선택에 따른 보수는 다음 표와 같다. (단, 표에서 각 셀 좌 하단의 숫자는 기업 A의 보수를, 우 상단의 숫자는 기업 B의 보수를 나타낸다.)
\begin{center}
\renewcommand{\arraystretch}{1.2}
\begin{tabular}{cc|c|c|}
\multicolumn{2}{c}{} & \multicolumn{2}{c}{\uline{B}} \\
\multicolumn{2}{c}{} & \multicolumn{1}{c}{L} & \multicolumn{1}{c}{H} \\ \cline{3-4}
\multirow{2}{*}{\uline{A}} & L & \parbox[c]{1.4cm}{{\raggedleft 1\par}\vspace{4pt}{\raggedright 1\par}} & \parbox[c]{1.4cm}{{\raggedleft 0\par}\vspace{4pt}{\raggedright 0\par}} \\ \cline{3-4}
 & H & \parbox[c]{1.4cm}{{\raggedleft 0\par}\vspace{4pt}{\raggedright 0\par}} & \parbox[c]{1.4cm}{{\raggedleft 0\par}\vspace{4pt}{\raggedright 0\par}} \\ \cline{3-4}
\end{tabular}
\end{center}
다음 설명 중 옳은 것을 모두 고른 것은?
\begin{framedtext}
ㄱ. 기업 A의 생산 전략 H는 우월전략이다. \par
ㄴ. 기업 A와 B 모두 생산 전략 L을 선택하는 것은 내쉬균형이다. \par
ㄷ. 기업 A와 B 모두 생산 전략 H를 선택하는 것은 내쉬균형이다.
\end{framedtext}
\begin{choices}
  \item ㄱ
  \item ㄴ
  \item ㄷ
  \item ㄱ, ㄴ
  \item ㄴ, ㄷ
\end{choices}
\end{question}
\answer{5}
